% ═══════════════════════════════════════════════════════════════════════════
%  فصل ۵ — ایران در سده‌های میانه: اعراب، ترکان و مغول‌ها
% ═══════════════════════════════════════════════════════════════════════════
\chapter{ایران در سده‌های میانه: اعراب، ترکان و مغول‌ها}
\label{ch:medieval-iran}

\begin{tcolorbox}[mybox, title={چکیده‌ی فصل}]
این فصل به بررسی نمونه‌های نجات‌طلبی و مداخله‌ی خارجی در ایران از فتح عرب تا پایان دوره‌ی تیموری (قرن هفتم تا نهم هجری / سیزدهم تا پانزدهم میلادی) می‌پردازد. تمرکز بر سه رویداد/فرایند اصلی است: تعامل حکومت‌های محلی ایرانی با خلافت عربی، حمله‌ی مغول و زمینه‌های داخلی آن، و الگوی تکرارشونده‌ی دعوت از ترکان و مغولان توسط نخبگان محلی. تحلیل بر مبنای مدل شش‌وجهی و با تکیه بر منابع فارسی، عربی و مغولی صورت می‌گیرد.
\end{tcolorbox}

% ─────────────────────────────────────────────────────────────────────────
\section{دوره‌ی خلافت و حکومت‌های نیمه‌مستقل ایرانی (قرن ۲–۴ هجری)}
\label{sec:ch05-caliphate}
% ─────────────────────────────────────────────────────────────────────────

\subsection{نهضت‌های ایرانی و مسئله‌ی «استقلال از درون»}
\label{subsec:ch05-iranian-movements}

پس از فتح عرب، ایرانیان دو استراتژی اصلی برای بازیابی هویت و قدرت سیاسی در پیش گرفتند:

\begin{enumerate}[nosep, label=\textbf{\arabic*.}]
  \item \textbf{مقاومت مستقیم:} شورش‌های مسلحانه مانند شورش \histref{بابک خرمدین}{۲۰۰–۲۲۲ ق / ۸۱۶–۸۳۷ م} در آذربایجان و شورش‌های مختلف در خراسان و طبرستان.%
  \footnote{نفیسی، سعید. \textit{بابک خرمدین}. اساطیر، ۱۳۸۴.}
  \item \textbf{نفوذ از درون:} حکومت‌های ایرانی‌تبار (طاهریان، صفاریان، سامانیان) که در چارچوب خلافت اما با استقلال عملی فراوان عمل می‌کردند.
\end{enumerate}

نکته‌ی مهم آن است که هیچ‌یک از این نهضت‌ها از نیروی خارجی «غیرمسلمان» یاری نخواستند. \thinker{عبدالحسین زرین‌کوب} (\lr{Abdolhossein Zarrinkoob}) در «دو قرن سکوت» استدلال می‌کند که ایرانیان پس از اسلام، «مقاومت فرهنگی» را بر «استقلال سیاسی مطلق» ترجیح دادند.%
\footnote{زرین‌کوب، عبدالحسین. \textit{دو قرن سکوت}. سخن، ۱۳۳۰، ص ۱۵--۳۵.}

\begin{histQuote}{زرین‌کوب درباره‌ی استراتژی ایرانیان}
«ایرانیان به‌جای آنکه اسلام را برانند، اسلام را ایرانی کردند. نه با شمشیر بیگانه بلکه با قلم و فرهنگ خود ... این شاید بزرگ‌ترین «انتقام» تاریخی بود: فاتح را از درون تغییر دادن.»
\end{histQuote}

\subsection{یعقوب لیث و الگوی «استقلال بدون استمداد»}
\label{subsec:ch05-yaqub}

\histref{یعقوب لیث صفاری}{۲۴۷–۲۶۵ ق / ۸۶۱–۸۷۹ م} نمونه‌ی بارز رهبری است که بدون هیچ اتکایی به نیروی خارجی، از قدرت‌های محلی برخاسته و علیه خلافت عباسی قیام کرد. پاسخ معروف او به فرستاده‌ی خلیفه — «این تیغ من و این کفش من! شما را به خلافت چه کار!» — بیان‌گر ایدئولوژی استقلال‌طلبانه‌ای بود که نیازی به «ناجی خارجی» نمی‌دید.%
\footnote{تاریخ سیستان، ص ۲۰۰--۲۱۵.}

\begin{tcolorbox}[defbox, title={نکته‌ی تحلیلی: چرا ایرانیان دوره‌ی اسلامی کمتر از بیگانه یاری خواستند؟}]
دلایل احتمالی:
\begin{enumerate}[nosep, label=\textbf{\arabic*.}]
  \item \textbf{فقدان «بیگانه‌ی قابل‌دسترس»:} بر خلاف دوره‌ی ساسانی که بیزانس همسایه‌ی مستقیم بود، پس از فتح عرب، تمام همسایگان ایران مسلمان بودند (یا مانند چین، بسیار دور).
  \item \textbf{هم‌کیشی:} اسلام به‌عنوان دین مشترک، مرز «خودی/بیگانه» را بازتعریف کرده بود.
  \item \textbf{حافظه‌ی تاریخی:} تجربه‌ی تلخ فتح عرب، ایرانیان را نسبت به «نجات‌دهنده‌ی خارجی» بی‌اعتماد کرده بود.
  \item \textbf{استراتژی نفوذ از درون:} ایرانیان دریافتند که نفوذ فرهنگی و بوروکراتیک در دستگاه خلافت مؤثرتر از مقاومت نظامی است.
\end{enumerate}
\end{tcolorbox}

% ─────────────────────────────────────────────────────────────────────────
\section{ورود ترکان: از مملوک تا سلطان}
\label{sec:ch05-turks}
% ─────────────────────────────────────────────────────────────────────────

\subsection{ترکان غزنوی و سلجوقی: دعوت‌شده یا مهاجم؟}
\label{subsec:ch05-ghaznavid-seljuq}

ورود ترکان به صحنه‌ی سیاسی ایران فرایندی تدریجی بود. خلفای عباسی از قرن سوم هجری، غلامان ترک را به‌عنوان سربازان وفادار وارد ارتش خود کردند — \concept{مملوک‌سازی} (\lr{Mamlukization}). این فرایند عملاً یک \textbf{نجات‌طلبی معکوس} بود: خلیفه از نیروی خارجی (ترکان) برای تحکیم قدرت خود در برابر نخبگان عرب و ایرانی یاری خواست.%
\footnote{\lr{Gordon, Matthew S. \textit{The Breaking of a Thousand Swords: A History of the Turkish Military of Samarra}. SUNY Press, 2001.}}

ورود \concept{سلجوقیان} به ایران (قرن پنجم هجری) نمونه‌ی پیچیده‌تری بود. \thinker{کلیفورد ادموند باسورث} (\lr{C.E. Bosworth}) نشان داده که برخی حکام محلی ایرانی — مانند امرای بلخ — عملاً از سلجوقیان علیه غزنویان یاری خواستند.%
\footnote{\lr{Bosworth, C.E. "The Political and Dynastic History of the Iranian World." In \textit{Cambridge History of Iran}, vol.~5, ed. Frye. Cambridge UP, 1968, pp.~1--202.}}

\begin{longtable}{
  >{\raggedleft\arraybackslash\bfseries}p{2.5cm}
  >{\raggedleft\arraybackslash}p{3cm}
  >{\raggedleft\arraybackslash}p{3cm}
  >{\raggedleft\arraybackslash}p{2.5cm}
  >{\raggedleft\arraybackslash}p{2cm}}
\caption{نمونه‌های دعوت/استقبال از ترکان در ایران}\label{tab:ch05-turk-invitation}\\
\toprule
\rowcolor{PrimaryDeep}
\textcolor{white}{\textbf{مورد}} &
\textcolor{white}{\textbf{دعوت‌کننده}} &
\textcolor{white}{\textbf{مداخله‌گر}} &
\textcolor{white}{\textbf{هدف}} &
\textcolor{white}{\textbf{نتیجه}} \\
\midrule
\endfirsthead
\toprule
\rowcolor{PrimaryDeep}
\textcolor{white}{\textbf{مورد}} &
\textcolor{white}{\textbf{دعوت‌کننده}} &
\textcolor{white}{\textbf{مداخله‌گر}} &
\textcolor{white}{\textbf{هدف}} &
\textcolor{white}{\textbf{نتیجه}} \\
\midrule
\endhead
\bottomrule
\endlastfoot

مملوک‌سازی عباسی &
خلفای عباسی &
غلامان ترک &
تحکیم قدرت خلیفه &
غلبه‌ی ترکان بر خلیفه%
\footnote{\lr{Kennedy, Hugh. \textit{The Prophet and the Age of the Caliphates}. 3rd ed. Routledge, 2016, pp.~185--210.}} \\
\rowcolor{NeutralLight}
ورود سلجوقیان &
برخی امرای خراسان &
سلجوقیان &
رهایی از غزنویان &
سلطه‌ی سلجوقی \\
دعوت از خوارزمشاهیان &
نخبگان محلی بلخ &
خوارزمشاه &
رهایی از سلجوقیان &
بی‌ثباتی بیشتر \\
\rowcolor{NeutralLight}
اسماعیلیان و مغول &
(نه دعوت مستقیم) &
مغول‌ها &
نابودی دشمنان &
فاجعه‌ی تمدنی \\

\end{longtable}

% ─────────────────────────────────────────────────────────────────────────
\section{حمله‌ی مغول: بزرگ‌ترین فاجعه‌ی تاریخ ایران}
\label{sec:ch05-mongol}
% ─────────────────────────────────────────────────────────────────────────

\subsection{زمینه‌ها و پرسش کلیدی}
\label{subsec:ch05-mongol-background}

حمله‌ی \histref{مغول‌ها}{۶۱۶–۶۵۴ ق / ۱۲۱۹–۱۲۵۶ م} به ایران فاجعه‌ی بزرگ تمدنی بود: تخمین‌ها حاکی از کشته‌شدن ۵ تا ۱۵ میلیون نفر (از جمعیت تقریبی ۲۰ میلیونی)، ویرانی زیرساخت‌های آبیاری، و نابودی کتابخانه‌ها و مراکز علمی است.%
\footnote{\lr{Morgan, David. \textit{The Mongols}. 2nd ed. Blackwell, 2007, pp.~73--82.}}

پرسش کلیدی: \textbf{آیا بخشی از ایرانیان از مغول‌ها دعوت کردند یا با آنها همکاری نمودند؟}

\subsection{شواهد «همکاری» با مغول‌ها}
\label{subsec:ch05-mongol-collaboration}

\subsubsection{خواجه نصیرالدین طوسی و نابودی خلافت عباسی}

\histref{خواجه نصیرالدین طوسی}{۵۹۷–۶۷۲ ق} مهم‌ترین چهره‌ی ایرانی است که به «همکاری» با مغول‌ها شهرت یافته. او پس از سال‌ها اسارت نزد اسماعیلیان، به خدمت \enemy{هولاکو} درآمد و نقش مشاوری در فتح بغداد و نابودی خلافت عباسی (۶۵۶ ق / ۱۲۵۸ م) ایفا کرد.%
\footnote{مدرسی رضوی، محمدتقی. \textit{احوال و آثار خواجه نصیرالدین طوسی}. بنیاد فرهنگ ایران، ۱۳۵۴، ص ۱۰۰--۱۳۵.}

\begin{tcolorbox}[casebox, title={تحلیل: خواجه نصیر — خائن یا عمل‌گرا؟}]
تفسیرهای مختلف:
\begin{itemize}[nosep, rightmargin=1em]
  \item \textbf{تفسیر سنتی شیعی:} خواجه نصیر «ناجی تشیع» بود که از قدرت مغول برای نابودی دشمن (خلیفه‌ی سنّی) و نجات شیعیان استفاده کرد.
  \item \textbf{تفسیر سنتی سنّی:} خواجه نصیر «خائنی» بود که به بیگانه‌ی کافر خدمت کرد و ویرانی بغداد بر عهده‌ی اوست.%
  \footnote{ابن‌تیمیه، تقی‌الدین. \textit{منهاج السنة النبویة}. تحقیق محمد رشاد سالم. مؤسسة قرطبه، ۱۹۸۶، ج ۳.}
  \item \textbf{تفسیر مدرن:} خواجه نصیر دانشمندی بود که در شرایط بحرانی، از فرصت بهره برد تا نهادهای علمی (رصدخانه‌ی مراغه) را حفظ کند — نه ناجی و نه خائن، بلکه \concept{عمل‌گرای بقا}.%
  \footnote{\lr{Dabashi, Hamid. "Khwajah Nasir al-Din Tusi." In \textit{History of Islamic Philosophy}, ed. Nasr and Leaman. Routledge, 1996, pp.~527--584.}}
\end{itemize}
\end{tcolorbox}

\subsubsection{حکام محلی و تسلیم به مغول‌ها}

\begin{longtable}{
  >{\raggedleft\arraybackslash\bfseries}p{2.5cm}
  >{\raggedleft\arraybackslash}p{3cm}
  >{\raggedleft\arraybackslash}p{3.5cm}
  >{\raggedleft\arraybackslash}p{3.5cm}}
\caption{واکنش حکام محلی ایران به حمله‌ی مغول}\label{tab:ch05-mongol-reactions}\\
\toprule
\rowcolor{PrimaryDeep}
\textcolor{white}{\textbf{حاکم/منطقه}} &
\textcolor{white}{\textbf{واکنش}} &
\textcolor{white}{\textbf{نتیجه‌ی فوری}} &
\textcolor{white}{\textbf{تحلیل}} \\
\midrule
\endfirsthead
\toprule
\rowcolor{PrimaryDeep}
\textcolor{white}{\textbf{حاکم/منطقه}} &
\textcolor{white}{\textbf{واکنش}} &
\textcolor{white}{\textbf{نتیجه‌ی فوری}} &
\textcolor{white}{\textbf{تحلیل}} \\
\midrule
\endhead
\bottomrule
\endlastfoot

جلال‌الدین خوارزمشاه &
مقاومت مسلحانه‌ی سرسختانه &
شکست نهایی و مرگ (۶۲۸ ق) &
تنها مقاومت سازمان‌یافته‌ی معنادار%
\footnote{نسوی، شهاب‌الدین. \textit{سیرت جلال‌الدین مینکبرنی}. ترجمه‌ی محمدعلی ناصح. علمی و فرهنگی، ۱۳۶۵.} \\
\rowcolor{NeutralLight}
اتابکان فارس &
تسلیم فوری و پرداخت خراج &
حفظ حکومت محلی &
عمل‌گرایی بقا \\
اسماعیلیان (حشاشین) &
مقاومت اولیه؛ سپس تسلیم &
نابودی کامل قلاع &
شکست مقاومت%
\footnote{\lr{Daftary, Farhad. \textit{The Ismāʿīlīs: Their History and Doctrines}. 2nd ed. Cambridge UP, 2007, pp.~355--398.}} \\
\rowcolor{NeutralLight}
خلیفه‌ی عباسی &
عدم آمادگی؛ تلاش دیرهنگام &
سقوط بغداد و قتل‌عام &
شکست فاجعه‌بار \\
بزرگان هرات &
مقاومت؛ سپس تسلیم؛ سپس شورش &
ویرانی کامل شهر در شورش دوم &
فاجعه%
\footnote{جوینی، عطاملک. \textit{تاریخ جهانگشای}. تصحیح محمد قزوینی. لیدن، ۱۹۱۲، ج ۱، ص ۱۲۵--۱۵۰.} \\

\end{longtable}

\subsection{نقش عطاملک جوینی و خاندان‌های ایرانی در دستگاه مغول}
\label{subsec:ch05-juvaini}

\histref{عطاملک جوینی}{۶۲۳–۶۸۱ ق}، نویسنده‌ی «تاریخ جهانگشای»، خود نمونه‌ای از نخبگان ایرانی است که در دستگاه مغول خدمت کردند. خاندان جوینی نسل‌ها در خدمت خوارزمشاهیان بود و سپس به‌سرعت به خدمت مغول‌ها درآمد.%
\footnote{\lr{Boyle, J.A. "The Death of the Last ʿAbbāsid Caliph." \textit{Journal of Semitic Studies} 6 (1961): 145--161.}}

\thinker{جان اندرو بویل} (\lr{J.A. Boyle}) این الگو را «پراگماتیسم بوروکراتیک» می‌نامد: نخبگان ایرانی با هر قدرت فاتحی همکاری می‌کردند — نه از سر نجات‌طلبی، بلکه از سر بقا و حفظ نهادهای اداری.%
\footnote{\lr{Boyle, J.A. "Dynastic and Political History of the Īl-Khāns." In \textit{Cambridge History of Iran}, vol.~5, ed. Boyle. Cambridge UP, 1968, pp.~303--421.}}

\begin{tcolorbox}[wavebox, title={الگوی «نفوذ از درون»: از ساسانیان تا مغول‌ها}]
الگوی تکرارشونده‌ای در تاریخ ایران وجود دارد: \textbf{نخبگان ایرانی به‌جای مقاومت نظامی یا نجات‌طلبی، سعی می‌کنند فاتح را از درون «ایرانی» کنند.}

\begin{itemize}[nosep, rightmargin=1em]
  \item ایرانیان دوره‌ی اسلامی خلافت را «ایرانی» کردند (وزرای برمکی، آل‌بویه و غیره).
  \item ایرانیان دوره‌ی مغول، ایلخانان را «اسلامی-ایرانی» کردند (غازان‌خان و اسلام‌آوردنش).
  \item این الگو در دوره‌ی تیموری و حتی صفوی نیز تکرار شد.
\end{itemize}

\thinker{برت فراگنر} (\lr{Bert Fragner}) این پدیده را «جهان‌شمول‌سازی فارسی» (\lr{Persianate universalism}) نامیده و آن را مهم‌ترین استراتژی بقای فرهنگ ایرانی دانسته است.%
\footnote{\lr{Fragner, Bert G. "Die 'Persophonie': Regionalität, Identität und Sprachkontakt in der Geschichte Asiens." ANOR 5. Berlin, 1999.}}
\end{tcolorbox}

% ─────────────────────────────────────────────────────────────────────────
\section{تیمور و نجات‌طلبی‌های دوره‌ی تیموری}
\label{sec:ch05-timur}
% ─────────────────────────────────────────────────────────────────────────

\subsection{حمله‌ی تیمور: تکرار فاجعه‌ی مغول}
\label{subsec:ch05-timur-invasion}

\histref{تیمور لنگ}{۷۳۶–۸۰۷ ق / ۱۳۳۶–۱۴۰۵ م} با ادعای «احیای شاهنشاهی چنگیز» به ایران حمله کرد. برخلاف مغول‌ها، تیمور مسلمان بود و از مشروعیت مذهبی برای فتوحاتش استفاده می‌کرد.%
\footnote{\lr{Manz, Beatrice Forbes. \textit{The Rise and Rule of Tamerlane}. Cambridge UP, 1989.}}

\subsection{دعوت‌های مستند از تیمور}
\label{subsec:ch05-timur-invitation}

برخی شواهد تاریخی نشان‌دهنده‌ی آن است که بخشی از نخبگان ایرانی عملاً از تیمور «دعوت» کردند:

\begin{enumerate}[nosep, label=\textbf{\arabic*.}]
  \item \textbf{علمای شیراز:} بر اساس روایت شرف‌الدین یزدی در «ظفرنامه»، برخی علمای شیراز نامه‌هایی به تیمور نوشتند و از ستم \enemy{شاه منصور} (آخرین مظفری) شکایت کردند و ورود تیمور را خواستند.%
  \footnote{شرف‌الدین علی یزدی. \textit{ظفرنامه}. تصحیح محمد عباسی. امیرکبیر، ۱۳۳۶، ج ۲، ص ۵۵--۷۰.}

  \item \textbf{سادات و اشراف اصفهان:} برخی سادات اصفهانی نیز تسلیم تیمور شدند و از او خواستند شهر را از حکام محلی «نجات» دهد — اما پس از شورش بخشی از مردم، تیمور دستور قتل‌عام داد و حدود ۷۰٬۰۰۰ نفر کشته شدند.%
  \footnote{\lr{Manz, 1989, pp.~71--74.}}
\end{enumerate}

\begin{tcolorbox}[enemybox, title={فاجعه‌ی اصفهان: درس تاریخی نجات‌طلبی}]
قتل‌عام اصفهان توسط تیمور (۷۸۹ ق / ۱۳۸۷ م) یکی از تلخ‌ترین نمونه‌های تاریخی پیامدهای نجات‌طلبی است:
\begin{enumerate}[nosep, label=\textbf{\arabic*.}]
  \item نخبگانی که از تیمور «دعوت» کرده بودند، تسلیم او شدند.
  \item بخشی از مردم عادی — که نه نجات‌طلب بودند و نه از دعوت خبر داشتند — علیه سربازان مغول شوریدند.
  \item تیمور در واکنش، فرمان قتل‌عام عمومی داد.
  \item نتیجه: \textbf{هم دعوت‌کنندگان و هم غیردعوت‌کنندگان قربانی شدند} — تجسم عینی ضرب‌المثل «مست گیرند، در شهر هر آنکه هست گیرند.»
\end{enumerate}
\end{tcolorbox}

% ─────────────────────────────────────────────────────────────────────────
\section{نمودار تحلیلی: الگوی نجات‌طلبی در قرون میانه‌ی ایران}
\label{sec:ch05-analysis-chart}
% ─────────────────────────────────────────────────────────────────────────

\begin{figure}[htbp]
\centering
\begin{tikzpicture}[
  every node/.style={font=\small},
  block/.style={
    draw=PrimaryMid, fill=BgCool, rounded corners=4pt,
    text width=3.5cm, align=center, minimum height=1.5cm,
    font=\small
  },
  arrow/.style={->, >=Stealth, thick, PrimaryMid},
  darrow/.style={->, >=Stealth, thick, AccentRed, dashed},
]

% ستون چپ: عوامل
\node[block, fill=BgWarm] (a1) at (0,6) {\rl{ضعف حکومت مرکزی}};
\node[block, fill=BgWarm] (a2) at (0,3.5) {\rl{نارضایتی نخبگان محلی}};
\node[block, fill=BgWarm] (a3) at (0,1) {\rl{تعدد مدعیان قدرت}};

% ستون وسط: فرایند
\node[block, fill=AccentGold!20, text width=4cm, minimum height=2cm] (process) at (6,3.5)
  {\rl{\textbf{نجات‌طلبی / همکاری}\\%
   با قدرت خارجی\\%
   (ترکان، مغول، تیمور)}};

% ستون راست: نتایج
\node[block, fill=BgSuccess] (r1) at (12,6) {\rl{حفظ جان و مقام فردی\\%
   (کوتاه‌مدت)}};
\node[block, fill=BgFail] (r2) at (12,3.5) {\rl{ویرانی و قتل‌عام\\%
   (میان‌مدت)}};
\node[block, fill=AccentPurple!15] (r3) at (12,1) {\rl{«ایرانی‌سازی» فاتح\\%
   (بلندمدت)}};

% فلش‌ها
\draw[arrow] (a1) -- (process);
\draw[arrow] (a2) -- (process);
\draw[arrow] (a3) -- (process);
\draw[arrow] (process) -- (r1);
\draw[darrow] (process) -- (r2);
\draw[arrow] (process) -- (r3);

\end{tikzpicture}
\caption{الگوی نجات‌طلبی و نتایج آن در قرون میانه‌ی ایران}
\label{fig:ch05-pattern}
\end{figure}

% ─────────────────────────────────────────────────────────────────────────
\section{مقایسه‌ی تطبیقی: مرزنشینان و غیرمرزنشینان}
\label{sec:ch05-border-interior}
% ─────────────────────────────────────────────────────────────────────────

آیا مرزنشینان بیشتر از ساکنان مرکز ایران به نجات‌طلبی متمایل بودند؟ شواهد تاریخی نشان می‌دهد:

\begin{table}[htbp]
\centering
\caption{مقایسه‌ی رفتار مرزنشینان و غیرمرزنشینان در مواجهه با مهاجمان}\label{tab:ch05-border}
\begin{adjustbox}{max width=\linewidth}
\begin{tabular}{
  >{\raggedleft\arraybackslash\bfseries}p{3cm}
  >{\raggedleft\arraybackslash}p{4cm}
  >{\raggedleft\arraybackslash}p{4cm}
  >{\raggedleft\arraybackslash}p{3cm}}
\toprule
\rowcolor{PrimaryDeep}
\textcolor{white}{ویژگی} &
\textcolor{white}{مرزنشینان} &
\textcolor{white}{ساکنان مرکز} &
\textcolor{white}{توضیح} \\
\midrule

تماس با بیگانه &
مکرر و مستقیم &
نادر و غیرمستقیم &
مرزنشینان با مهاجمان آشنایی بیشتری داشتند \\
\rowcolor{NeutralLight}
وفاداری به مرکز &
مشروط و متغیر &
نسبتاً قوی‌تر &
وفاداری مرزنشینان تابع حمایت مرکز بود \\
تمایل به تسلیم &
بالاتر (واقع‌گرایی) &
پایین‌تر (ایدئالیسم) &
مرزنشینان نخستین قربانیان جنگ بودند \\
\rowcolor{NeutralLight}
نقش اقتصادی &
تجارت مرزی، راه‌های بازرگانی &
کشاورزی، صنعت &
مرزنشینان از صلح با بیگانه سود اقتصادی می‌بردند \\
تنوع قومی-مذهبی &
بالاتر &
پایین‌تر &
تنوع قومی وفاداری انحصاری به یک مرکز را ضعیف می‌کرد \\

\bottomrule
\end{tabular}
\end{adjustbox}
\end{table}

% ─────────────────────────────────────────────────────────────────────────
\section{جمع‌بندی فصل}
\label{sec:ch05-summary}
% ─────────────────────────────────────────────────────────────────────────

\begin{tcolorbox}[wavebox, title={یافته‌های کلیدی فصل ۵}]
\begin{enumerate}[nosep, label=\textbf{\arabic*.}]
  \item در دوره‌ی اسلامی اولیه (قرن ۲–۴ ق)، ایرانیان به‌ندرت از «بیگانه‌ی غیرمسلمان» یاری خواستند. استراتژی غالب، \concept{مقاومت فرهنگی} و \concept{نفوذ از درون} بود.

  \item ورود ترکان به‌عنوان نیروی نظامی، خود نوعی «نجات‌طلبی معکوس» توسط خلفای عباسی بود که در نهایت به غلبه‌ی ترکان بر خلیفه انجامید — تأیید الگوی «مست گیرند...».

  \item در مواجهه با مغول و تیمور، الگوی غالب «تمکین عملی» بود نه «نجات‌طلبی». موارد معدود «دعوت» (مانند علمای شیراز) نتایج فاجعه‌باری داشت (قتل‌عام اصفهان).

  \item نخبگان ایرانی (مانند خواجه نصیر و جوینی) استراتژی «ایرانی‌سازی فاتح» را در پیش گرفتند — استراتژی‌ای که در بلندمدت مؤثرتر از نجات‌طلبی بود اما هزینه‌های اخلاقی سنگینی داشت (مشارکت در قتل‌عام‌ها).

  \item مرزنشینان نسبت به ساکنان مرکز، تمایل بیشتری به تسلیم و همکاری با مهاجم داشتند — نه لزوماً از سر «خیانت» بلکه از سر \concept{واقع‌گرایی بقایی} (\lr{survival pragmatism}).

  \item الگوی «نفوذ از درون» استراتژی ویژه‌ی ایرانی بود که در هیچ تمدن دیگری (مثلاً بین‌النهرین عربی یا هند) به این شکل مشاهده نمی‌شود. \thinker{مارشال هاجسن} (\lr{Marshall Hodgson}) این پدیده را بخشی از «جهان فارسی‌زبان» (\lr{Persianate World}) و قدرت نرم تمدن ایرانی دانسته است.%
  \footnote{\lr{Hodgson, Marshall G.S. \textit{The Venture of Islam}. U of Chicago P, 1974, vol.~2, pp.~293--314.}}
\end{enumerate}
\end{tcolorbox}

\cleardoublepage